\documentclass[11pt,a4paper]{article}
\usepackage[T1]{fontenc}
\usepackage[utf8]{inputenc}
\usepackage{amsmath,amssymb,amsthm}
\usepackage[margin=1in]{geometry}
\usepackage[colorlinks=true,linkcolor=blue,urlcolor=blue]{hyperref}
\theoremstyle{plain}
\newtheorem{theorem}{Theorem}
\newtheorem{lemma}[theorem]{Lemma}
\newtheorem{proposition}[theorem]{Proposition}
\newtheorem{corollary}[theorem]{Corollary}
\theoremstyle{definition}
\newtheorem{remark}[theorem]{Remark}

\title{The empirical recurrence for OEIS A202893, proved by transfer matrix}
\author{Adrian Perez Fontelles\\ \small Independent researcher}
\date{3 September 2026}
\begin{document}
\maketitle

\begin{abstract}
OEIS A202893 counts arrays of width $5$ over $\{0,1\}$ containing no occurrence of certain short
patterns, read along the rows, the columns or the diagonals. The entry carries an empirical
recurrence of order $82$ contributed by R.~H.~Hardin. It is true, and it is decidable
rather than empirical: no pattern is longer than $3$ letters and no direction moves the row
index by more than one, so every occurrence lies inside $3$ consecutive array rows and the
count is a walk count on $S=211$ states.
\end{abstract}

\noindent\small 2020 Mathematics Subject Classification. 05A15, 68R15, 15A18.\normalsize

\section{The conjecture}

OEIS A202893 is ``Number of (n+2)X5 binary arrays avoiding patterns 001 and 010 in rows, columns and nw-to-se diagonals.''. It has offset $1$ and begins
\[
776, 3264, 17116, 81308, 369152, 1719383, 7994866, 36814964,\ \dots
\]
The entry states, as a conjecture contributed by its author and never marked settled:
\begin{quote}
Empirical: a(n) = 7*a(n-1) -20*a(n-2) +107*a(n-3) -398*a(n-4) +738*a(n-5) -2666*a(n-6) +6567*a(n-7) -5696*a(n-8) +23763*a(n-9) -38174*a(n-10) -28818*a(n-11) -128857*a(n-12) +91998*a(n-13) +497173*a(n-14) +980615*a(n-15) +2613*a(n-16) -2944178*a(n-17) -6692446*a(n-18) -3201580*a(n-19) +12324831*a(n-20) +31942303*a(n-21) +27028612*a(n-22) -33117909*a(n-23) -116041590*a(n-24) -124565023*a(n-25) +41055880*a(n-26) +313353070*a(n-27) +407261622*a(n-28) +54971014*a(n-29) -612525762*a(n-30) -981333878*a(n-31) -429317040*a(n-32) +851543798*a(n-33) +1753792532*a(n-34) +1184415382*a(n-35) -726958892*a(n-36) -2352756541*a(n-37) -2083099763*a(n-38) +90662131*a(n-39) +2318885002*a(n-40) +2613663583*a(n-41) +745650465*a(n-42) -1576402440*a(n-43) -2392059010*a(n-44) -1253899240*a(n-45) +596726635*a(n-46) +1575476433*a(n-47) +1177422401*a(n-48) +75440229*a(n-49) -705283171*a(n-50) -737641009*a(n-51) -273495008*a(n-52) +169039353*a(n-53) +312071492*a(n-54) +197384459*a(n-55) +18468996*a(n-56) -80302791*a(n-57) -82414262*a(n-58) -34856020*a(n-59) +5887432*a(n-60) +20946985*a(n-61) +15251291*a(n-62) +4417667*a(n-63) -2598284*a(n-64) -3750121*a(n-65) -2034951*a(n-66) -220817*a(n-67) +543926*a(n-68) +447516*a(n-69) +154787*a(n-70) -34294*a(n-71) -59123*a(n-72) -29828*a(n-73) -1394*a(n-74) +4423*a(n-75) +3179*a(n-76) +350*a(n-77) -182*a(n-78) -169*a(n-79) -23*a(n-80) +5*a(n-81) +4*a(n-82).
\end{quote}
As of the ``Last modified'' line on the live entry (Jul 19 2023 10:17:05, revision 8) this is still
recorded as empirical, and nothing on the entry records it as proved.

\section{What is being counted}

The array has $5$ cells across, with entries in $\{0,1\}$. An OCCURRENCE of a pattern in a
direction is a run of consecutive cells taken in that direction whose entries match the pattern
in order; the array is counted when it has no occurrence at all. The entry names
\[
\begin{array}{ll}
\text{direction} & \text{forbidden patterns} \\[2pt]
\text{along a row, left to right} & \mathtt{001},\quad \mathtt{010} \\
\text{down a column} & \mathtt{001},\quad \mathtt{010} \\
\text{down a diagonal, nw to se} & \mathtt{001},\quad \mathtt{010} \\
\end{array}
\]


\section{Every occurrence lies in $3$ consecutive rows}

\begin{lemma}\label{lem:win}
Each of the directions above changes the row index by $0$ or $1$ from one cell of a run to the
next, and no pattern is longer than $3$ letters. Hence every occurrence lies inside $3$
consecutive array rows, and an occurrence whose last cell is in row $i$ lies inside rows
$i-2,\dots,i$.
\end{lemma}

\begin{proof}
Immediate from the offsets listed: the row step is $0$ for a row and $1$ for the other three.
\end{proof}

So the state is the window of the last $2$ array rows, a row still to come being written
$\bot$. A step appends a row and rejects it exactly when some occurrence ends in it or lies
inside it; an occurrence needing a row that is $\bot$ cannot exist, which is the boundary rule
and is not assumed but read off the window. The walk starts at the all-$\bot$ window, so a walk
of no steps is the empty array and a walk of $R$ steps is an array of $R$ rows. With $M$ the
adjacency matrix, $\iota$ the indicator of the starting window and $\tau=\mathbf{1}$,
\[
a(n)\;=\;\iota^{\!\top}M^{\,j}\tau ,\qquad S=211.
\]
The walk index $j$ and the entry's own $n$ differ by a constant, read off by matching the
published terms rather than assumed. In particular $a$ is $C$-finite.

\section{The proof}

\begin{lemma}\label{lem:crit}
Let $q(t)=t^{82}-\sum_i c_it^{82-i}$ be the characteristic polynomial of the
conjectured recurrence. The residual $a(n)-\sum_ic_ia(n-i)$ equals
$\iota^{\!\top}M^{\,j}q(M)\tau$ for the corresponding walk index $j$, so the recurrence
holds from some point on if and only if $\iota^{\!\top}M^{j}q(M)\tau=0$ for all large $j$,
and it is enough to examine $j<S$.
\end{lemma}

\begin{proof}
Factor $M^{j}$ out of $M^{j+82}-\sum_ic_iM^{j+82-i}$. For the bound, $M^{S}$
is an integer combination of $I,M,\dots,M^{S-1}$ by Cayley--Hamilton, so the numbers
$u_j=\iota^{\!\top}M^{j}q(M)\tau$ obey the monic recurrence given by the characteristic
polynomial of $M$; $S$ consecutive zeros force every later one, and the last nonzero $u_j$
pins the threshold exactly.
\end{proof}

\begin{theorem}
\[
a(n)\;=\;7\,a(n-1) - 20\,a(n-2) + 107\,a(n-3) - 398\,a(n-4) + 738\,a(n-5) - 2666\,a(n-6) + 6567\,a(n-7) - 5696\,a(n-8) + 23763\,a(n-9) - 38174\,a(n-10) - 28818\,a(n-11) - 128857\,a(n-12) + 91998\,a(n-13) + 497173\,a(n-14) + 980615\,a(n-15) + 2613\,a(n-16) - 2944178\,a(n-17) - 6692446\,a(n-18) - 3201580\,a(n-19) + 12324831\,a(n-20) + 31942303\,a(n-21) + 27028612\,a(n-22) - 33117909\,a(n-23) - 116041590\,a(n-24) - 124565023\,a(n-25) + 41055880\,a(n-26) + 313353070\,a(n-27) + 407261622\,a(n-28) + 54971014\,a(n-29) - 612525762\,a(n-30) - 981333878\,a(n-31) - 429317040\,a(n-32) + 851543798\,a(n-33) + 1753792532\,a(n-34) + 1184415382\,a(n-35) - 726958892\,a(n-36) - 2352756541\,a(n-37) - 2083099763\,a(n-38) + 90662131\,a(n-39) + 2318885002\,a(n-40) + 2613663583\,a(n-41) + 745650465\,a(n-42) - 1576402440\,a(n-43) - 2392059010\,a(n-44) - 1253899240\,a(n-45) + 596726635\,a(n-46) + 1575476433\,a(n-47) + 1177422401\,a(n-48) + 75440229\,a(n-49) - 705283171\,a(n-50) - 737641009\,a(n-51) - 273495008\,a(n-52) + 169039353\,a(n-53) + 312071492\,a(n-54) + 197384459\,a(n-55) + 18468996\,a(n-56) - 80302791\,a(n-57) - 82414262\,a(n-58) - 34856020\,a(n-59) + 5887432\,a(n-60) + 20946985\,a(n-61) + 15251291\,a(n-62) + 4417667\,a(n-63) - 2598284\,a(n-64) - 3750121\,a(n-65) - 2034951\,a(n-66) - 220817\,a(n-67) + 543926\,a(n-68) + 447516\,a(n-69) + 154787\,a(n-70) - 34294\,a(n-71) - 59123\,a(n-72) - 29828\,a(n-73) - 1394\,a(n-74) + 4423\,a(n-75) + 3179\,a(n-76) + 350\,a(n-77) - 182\,a(n-78) - 169\,a(n-79) - 23\,a(n-80) + 5\,a(n-81) + 4\,a(n-82)
\]
for every $n>82$.
\end{theorem}

\begin{proof}
The vector $w=q(M)\tau$ was formed by $82$ matrix--vector products in exact integer
arithmetic and $\iota^{\!\top}M^{j}w$ evaluated for $j=0,1,\dots$, again exactly; those
integers vanish from the index corresponding to $n=82+1$ onwards and the run continued
until the working vector was identically zero, which settles every larger $j$ at once.
\end{proof}

\section{Verification}

Four checks, each independent of the algebra above.

First, the model reproduces all $19$ terms the entry publishes, exactly, in integer
arithmetic.

Second, the count was recomputed for the smallest arrays straight from the definition: fill the
cells one at a time, in the orientation the entry's own name writes, and reject as soon as a
completed run matches a forbidden pattern. No transfer matrix, no window and no transposition
are involved, and the published terms came back.

Third, the conjectured recurrence was evaluated directly on the published terms, with no
matrices involved, and holds wherever the proved range and the data overlap.

Fourth, the annihilation test of Lemma~\ref{lem:crit} was carried out in exact integer
arithmetic on the whole vector rather than on a sampled prefix.

\begin{thebibliography}{9}
\bibitem{oeis} The OEIS Foundation, \emph{The On-Line Encyclopedia of Integer
Sequences}, \url{https://oeis.org}, 2026. Sequence A202893.
\bibitem{stanley} R.~P.~Stanley, \emph{Enumerative Combinatorics}, Volume 1, 2nd ed.,
Cambridge University Press, 2011. (Transfer-matrix method, Section 4.7.)
\bibitem{hu} J.~E.~Hopcroft, R.~Motwani and J.~D.~Ullman, \emph{Introduction to Automata
Theory, Languages, and Computation}, 3rd ed., Addison--Wesley, 2006.
\bibitem{horn} R.~A.~Horn and C.~R.~Johnson, \emph{Matrix Analysis}, 2nd ed., Cambridge
University Press, 2013.
\end{thebibliography}

\end{document}
